%% AgenticOS 2026 (SOSP) vision paper submission.
%% Double-blind: do not identify authors or the system's public identity.
%% Bibliography is hand-rolled (thebibliography), so no bibtex pass is needed.
%% Compile:  pdflatex main && pdflatex main   (twice, for refs/table)
%% Or drop into Overleaf (acmart is standard there; pick pdflatex compiler).
%%
%% Evidence for the two measured rows lives in evidence/bench.json
%% (reproducible via evidence/bench.py against a running coforge-router).

\documentclass[sigconf, review, anonymous]{acmart}
%% 'anonymous' hides author block for double-blind review.
%% 'review' adds line numbers for reviewers.
%% acmart already loads hyperref, graphicx, xcolor — do not re-load.

%% Venue — replaces the default "Conference'17, July 2017, Washington, DC" placeholder.
%% Workshop has no formal proceedings, so we set copyright to none and omit
%% ISBN/DOI/price (which would otherwise render as fake placeholder values).
\setcopyright{none}
\settopmatter{printacmref=false, printccs=false, printfolios=false}
\acmConference{2nd Workshop on Operating Systems Design for AI Agents (AgenticOS)}{September 29, 2026}{Prague, Czechia}
\acmYear{2026}
%% No formal proceedings → suppress fake ISBN/DOI/price placeholders.
\acmISBN{}
\acmDOI{}
\acmPrice{}

\usepackage{booktabs}
\usepackage{array}

\begin{document}

\title{An Agent Workspace as a Measuring Instrument for the OS It Lacks}

\author{Anonymous}
\affiliation{\institution{Anonymized for double-blind review}}

\begin{abstract}
AI agents are becoming long-lived services that plan, invoke tools, and
persist state, yet they are still hosted by application-layer workspaces
built on operating-system abstractions designed for static workloads. We
built and ran such a workspace---an open-source, multi-agent chat
application with persistent per-agent memory---and report where its
implementation hit walls that are missing OS abstractions. We enumerate
six workarounds, measure all six (three scaling, three correctness cliffs),
and map each onto an OS abstraction that would remove it. The contribution
is consumer-side evidence, not a new OS, for the set of abstractions the
agent era is starved for.
\end{abstract}

\maketitle

\section{Introduction}

Operating systems expose processes, threads, files, sockets, and resource
controllers as the units of execution, composition, and accounting. These
abstractions were designed for workloads that are static, predictable, and
short-lived relative to human attention. AI agents invert all three
properties: they are long-lived, semantically adaptive, and interact with
tools and environments in ways the runtime cannot predict. A recent workshop
series argues that this mismatch is not incidental but structural, and that
operating systems themselves must become \emph{agentic}~\cite{agenticos}.

Most of the evidence offered for this claim comes from the
\emph{producer} side of the OS interface: new kernel primitives, new
schedulers, new isolation models. This paper offers evidence from the
\emph{consumer} side. We built and operate a real, open-source agent
workspace---a multi-agent chat application in which a human routes work to
named agents via channel mentions, and each agent retains memory across
sessions. The application is small (a few hundred lines) and states its
non-goals explicitly. We found that at every hard problem, the limiting
factor was not the application's logic but the absence of an
operating-system abstraction that should plausibly exist.

This is a position paper, not a new system: the contribution is the mapping.
Each of the workspace's admitted workarounds corresponds to an OS-level
abstraction that would dissolve it.
We measure all six to show they are not cosmetic: three scaling walls
(latency vs concurrency, tokens vs history, stored rows vs turns) and three
correctness cliffs (isolation, delegation, capability routing), each failing
on 100\% of fault-injection trials. Each is technically re-implementable in user space, but
only by re-building, per application, the scheduling and managed-state
primitives that belong one layer down---and by getting their corner cases
right every time.

\section{The workspace}

The system is an open-source agent workspace (anonymized for review). A
human types into a single chat channel; messages beginning with
\texttt{@Name} are routed to the named agent. Each agent has a persona and a
private, append-only memory. On each turn the agent's full history is
replayed into the prompt, the model is invoked through a standard
OpenAI-compatible endpoint, and the reply is stored back. Memory is
isolated per agent: a fact told to one agent is not visible to another.

The implementation comprises a Node.js HTTP router, a SQLite store for
channel history and per-agent memory, and a minimal web frontend. It uses
the host Node runtime's built-in SQLite module; no external agent
framework, vector database, or daemon is involved. The system is deployed
for our own use.

What makes the system useful as evidence is not its capability but the
explicitness of its non-goals. By design, it has:

\begin{itemize}
  \item \textbf{No multi-agent concurrency.} Messages in a channel are
        processed through one serial queue. Two agents cannot act on the
        same state at the same time.
  \item \textbf{No memory compression.} The entire history is replayed
        into the prompt on every turn. There is no summarization or
        archival layer.
  \item \textbf{No inter-agent isolation.} All agents share one process
        and one database; there is no boundary preventing one agent from
        reading another's state.
  \item \textbf{No capability routing.} A mention is parsed by a regular
        expression; there is no matching of work to the agent best suited
        for it, and no delegation between agents.
\end{itemize}

These are not oversights. Each is a deliberate non-goal that kept the
implementation readable. We now argue that each is also a place where an
operating-system abstraction is conspicuously absent.

\section{Workarounds and the abstractions they want}

Table~\ref{tab:map} lists the six workarounds and the OS abstraction each
implies. They split into three scaling-pathology walls (growth curves) and
three correctness-cliff walls (invisible until the first adversarial event,
surfaced via fault-injection rates). We measure all six below in the
uncompressed configuration (the workspace ships a summarization switch that
bounds prompt-token growth; we disable it to surface the walls). Single
runs on a developer laptop, one OpenAI-compatible model (glm-5.1),
reproducible from the artifact (anonymized for review).

\begin{table}[t]
\caption{Each workaround maps to an OS abstraction that would remove it.}
\label{tab:map}
\small
\setlength{\tabcolsep}{4pt}
\begin{tabular}{p{0.46\linewidth} p{0.46\linewidth}}
\toprule
Workaround (what is hand-rolled) & OS abstraction wanted \\
\midrule
Serial queue (no concurrency) & Semantics-aware agent scheduling \\
Full-history replay (no compression) & Long-lived managed state \\
One process + one DB (no boundary) & Execute-only / attested isolation \\
JSON personas (hand-authored) & Skills as unit of composition \\
\texttt{@Name} regex (no matching) & Agent-facing interface / IPC \\
SQLite file (no managed object) & Agent state as OS object \\
\bottomrule
\end{tabular}
\end{table}

\subsection{Measured: the serial queue}

\subsection{Measured: six walls, two instruments}

The six walls are measured with two instruments: a growth curve for the
three scaling-pathology walls, and a fault-injection rate for the three
correctness cliffs. Table~\ref{tab:measured} summarizes; the pattern is
that every scaling wall grows linearly and every cliff fails on 100\% of
trials. Each wall is the absence of the abstraction Table~\ref{tab:map}
names: a degenerate scheduler (no global order basis), full-history replay
(no compression), append-only storage (no eviction), an advisory boundary
(no runtime mediation), no inter-agent call path (no composition), and
@name regex (no capability matching). The application cannot fix any of
them without re-implementing, in user space, the abstraction one layer
down---and it cannot do isolation correctly at all, being both the enforcer
and the enforced-upon (\S\ref{sec:abstractions}).

\begin{table}[t]
\caption{All six walls measured. Scaling walls: growth shape. Cliffs: fault-injection cross-rate.}
\label{tab:measured}
\small
\setlength{\tabcolsep}{4pt}
\begin{tabular}{p{0.30\linewidth} p{0.18\linewidth} p{0.40\linewidth}}
\toprule
Wall & Family & Result \\
\midrule
Serial queue & scaling & latency 3.3$\to$27.3\,s over $N$=1..8 (linear) \\
Prompt replay & scaling & tokens 39$\to$308 over 8 turns (linear) \\
Managed state & scaling & rows 2$\to$40 over 20 turns (linear, no eviction) \\
Isolation & cliff & boundary crossed 100/100 trials \\
Composition & cliff & delegation reached target 0/20 trials \\
Routing & cliff & mismatched task refused 0/20 trials \\
\bottomrule
\end{tabular}
\end{table}

\section{Why these are abstractions, not features}
\label{sec:abstractions}

A reviewer may object that the items in Table~\ref{tab:map} are missing
\emph{features}, not missing \emph{abstractions}. The distinction is the
contribution, so we give a falsifiable test: \emph{a requirement is an
abstraction, not a feature, when any application-layer implementation of it
is necessarily best-effort and cannot be verified correct in isolation.} A
feature lives inside one program's logic and can be made correct there; an
abstraction is a contract that must be guaranteed by the layer beneath, or
not at all.

Isolation is the clearest case (Table~\ref{tab:measured} measures it): any
boundary an application draws is advisory---a bug, a misconfigured tool, or
a prompt-injected agent steps over it, and the application has no vantage
point to detect the crossing. Only the runtime, mediating every memory
access, can enforce one that holds. The same structure---enforcer and
enforced-upon are the same program---repeats across the other rows:
concurrency safety, context compression, capability routing each cannot be
verified by the process whose own behavior it constrains.

This is also why the walls are not specific to our workspace. Any
application that lets agents persist, collaborate, or be selected by
capability occupies the same position---enforcer and enforced-upon at
once---and hits the same walls, because the walls are in the contract
between the application and its runtime, not in the application's logic.
The fact that a few-hundred-line application is forced to hand-roll all of
these, badly, is the evidence.

\section{Relation to the AgenticOS agenda}

Each row aligns with a topic the AgenticOS community has identified from the
producer side: agent-aware resource control~\cite{agentcgroup},
long-lived state for agent memory~\cite{agenticos}, execute-only and
attested-channel isolation~\cite{execonly,grimlock}, skills as composition
unit~\cite{skillos}, agent-facing interface redesign~\cite{rethos}. But the
1st-edition agenda is weighted roughly ten-to-two toward performance over
security: the three scaling walls we measure sit on the well-served side,
while the three correctness cliffs---isolation, composition, routing---occupy
the under-served ground. The asymmetry in our evidence mirrors the asymmetry
the field carries.

\section{What we are not claiming}
\label{sec:notclaiming}

We do not claim a new operating system. We do not claim that the
abstractions named in Table~\ref{tab:map} are the \emph{correct} ones, only
that they are \emph{wanted}. We do not claim the workspace is
representative of all agent applications; we claim that any application
with its shape---persistent, multi-agent, capability-routed---will hit the
same walls because the walls are in the runtime contract. We do not claim
these abstractions are easy; we claim their absence is expensive, and we
have the measurements to show it. The measurements are limited: single-run,
one model, small N (8--20--100), no seeds or CIs. They show growth
\emph{shape} (three scaling walls) and cross-\emph{rates} (three cliffs),
not magnitudes. The skills-composition probe is a weak proxy (no scalar
metric for composition exists); the routing probe uses an LLM judge.

\enlargethispage{\baselineskip}

\section{Conclusion}

An application-layer agent workspace that states its non-goals explicitly is
a measuring instrument for the abstractions its operating system lacks. We
built one, measured all six of its walls (three scaling, three cliffs), and mapped
six to OS-level abstractions the agent era needs. Consumer-side evidence of
this kind complements the producer-side primitives the community is
beginning to propose, toward a contract agent applications request rather
than re-implement.

{\footnotesize
\begin{thebibliography}{9}

\bibitem{agenticos}
AgenticOS Workshop Series. \emph{Operating Systems Design for AI Agents.}
SOSP 2026 (2nd edition); ASPLOS 2026 (1st edition).
\url{https://os-for-agent.github.io/}

\bibitem{agentcgroup}
Y.~Zheng, J.~Fan, Q.~Fu, et~al.
\emph{AgentCgroup: Understanding and Controlling OS Resources of AI
Agents.} AgenticOS 2026.

\bibitem{skillos}
L.~Chen, Z.~Wang, W.~Zheng, et~al.
\emph{Skills are the new Apps---Now It Is Time for Skill OS.}
AgenticOS 2026.

\bibitem{execonly}
R.~Tiwari, D.~Williams. \emph{Execute-Only Agents: Architectural Defense
Against Prompt Injection for AI Agents.} AgenticOS 2026.

\bibitem{grimlock}
Q.~Wu, W.~Zhang, G.~Fang, et~al. \emph{Grimlock: Guarding High-Agency Systems
with eBPF and Attested Channels.} AgenticOS 2026.

\bibitem{rethos}
Y.~Wang et~al. \emph{Rethinking OS Interfaces for LLM Agents.}
AgenticOS 2026.

\end{thebibliography}
}

\balance
\end{document}
